\section{Event-level transport}\label{app:E.25}

A \SI{100}{eV} electron recoil in silicon ends its cascade with about 27 electron--hole pairs, at \SI{3.75}{eV} per pair~\cite{rodrigues2021absolute}, and with the rest of its energy in phonons. What the detector reports is decided afterwards, in a readout window of a microsecond to a second. In that window the pairs drift, bind into excitons or get trapped, the gap energy they carry is released somewhere as phonons, and each sensor absorbs some share of each phonon population. \Cref{sec:18.5,sec:19.6} and \cref{chap:20} derive the physics of every step, and \cref{app:E.12} sets out the transport Monte Carlo that chains the steps. The chapters do not supply the rates, kernels and efficiencies of a stated detector, and the readout conclusions in the Summary of \cref{chap:20} stay conditional until these are computed. Every calculation below sits at the place on the truncation label $\mathfrak{T}$ (\cref{sec:8.5}) that \cref{app:E.12} gives the transport: lowest-order golden-rule rates, with $n_g = 2$ for phonon emission by carriers, $n_A = 2$ for three-phonon decay, $n_B = 0$, and incoherent propagation between events. \Cref{tab:E25-runs} collects grids, acceptance tests and costs.

\subsection{Charge and exciton kinetics}

The measured charge and the measured heat of an event both depend on what the pairs do during the window $t_{\mathrm{win}}$, and both follow from the joint distribution of the measured pair number $N_{\mathrm{meas}}$ and of the energy $E_{\mathrm{st}}(t_{\mathrm{win}})$ still stored at the end of the window. It is conditional on the recoil energy $E_R$, the recoil direction $\hat{\vect{r}}_R$, the applied field $\mathcal{E}$, the window and the temperature $T$. From it follow two tables that \cref{sec:19.5,sec:19.6} use: the phonon-yield ratio $Y_{\mathrm{ph}}(t_{\mathrm{win}})$ of nuclear to electron recoils evaluated at a finite window, for $t_{\mathrm{win}} \in \{\SI{1}{\micro s}, \SI{100}{\micro s}, \SI{10}{ms}, \SI{1}{s}\}$, and the same ratio at detector thresholds of 0.3, 1, 3 and \SI{16}{eV}. The grid is 20 logarithmic recoil energies from \SI{10}{eV} to \SI{10}{keV}, fields of 0, 1, 4 and \SI{40}{V/cm}, and the impurity census of the detector modeled (boron, phosphorus and the deep centers its growth certificate lists, at their stated concentrations).

That distribution comes from the rate equations of \cref{sec:18.5} for free pairs, free and bound excitons and, at high density, electron--hole droplets, coupled to drift; the luminescent share leaves the crystal as escaping energy $E_{\mathrm{esc}}$ (\cref{sec:19.5}). The field enters in two ways: pairs separated by drift faster than the binding time do not form excitons, and above $E_x^b/(ea_x) = \SI{3.5e4}{V/cm}$ excitons are field-ionized (\cref{sec:18.5}). Drift, Luke emission, trapping and Ramo induction are those of the carrier Monte Carlo of \cref{app:E.12}, with four additions: the trap occupations $n_j(t)$ carry over from event to event, trap energies are referred to the Fermi level, the electrostatic energy of the trapped charge is a separate share, and the drift during the prompt phase is kept. The exciton module is new code. The exciton and capture parameters come from \cref{app:E.24}.

Two measured numbers serve as acceptance tests. The silicon phonon-yield ratio of the CDMS~II detectors, $Y_{\mathrm{ph}} = 0.952^{+0.009}_{-0.007}$ (\SI{95}{\percent} confidence) at \SI{4}{V/cm}~\cite[Sect.~4]{agnese2018nuclear}, and the defect energy loss of $^{206}$Pb recoils in germanium, $(6.08 \pm 0.18)\,\%$~\cite{agnese2018energy}, must both be reproduced by one parameter set. The Hecht limit of \cref{app:E.12} is the numerical check. The calculation is meaningful only together with the collection efficiencies and the athermal emission rates specified below, since the measured yield mixes all three. A parameter set costs up to $10^{3}$ core-hours; the full table over windows, thresholds and the two materials $10^{3}$ to $10^{4}$.

\subsection{Athermal emission and space charge}

At millikelvin temperatures thermal emission from traps is frozen out, and two effects remain that a constant trapping length cannot describe. Trapped charge is released athermally, by phonons and by hot carriers, and the charge that stays trapped distorts the drift field. The calculation gives the emission rates $\varGamma^{\mathrm{em}}_j$ of each trap class $j$, in \si{s^{-1}}, over a stated exposure history, and the field distortion $\xi_{\mathrm{SC}} = \max_z|\Delta\mathcal{E}|/\mathcal{E}$ of \cref{sec:20.1}.

The field distortion follows from the one-dimensional Poisson equation for a sheet of trapped charge (\cref{eq:20.1-sheet}). In silicon $\xi_{\mathrm{SC}} = \SI{1}{\percent}$ at \SI{1}{V/cm} needs 6.3 to \SI{12.6e4}{cm^{-2}} trapped charges of one sign, which \SI{100}{eV} events at \SI{1}{Hz} on \SI{1}{cm^2} supply in 5 to 10 hours and dense calibration runs in minutes (\cref{sec:20.1}). The grid satisfies $\dd z \le d/200$ for a crystal of thickness $d$, and the event sequence runs until the distortion saturates or is reset by flashing with infrared light.

The emission rates need cross-sections that are not known at millikelvin. A neutral phosphorus donor in silicon is bound by \SI{45}{meV}, and phonons above \SI{10.9}{THz} live below a nanosecond and travel about a micrometer (\cref{sec:20.2}), so phonon-driven emission from neutral shallow centers happens only next to the track of a later event. The Luke phonons, a few meV, cannot ionize neutral centers, but can ionize the over-charged centers $\mathrm{D}^-$ and $\mathrm{A}^+$, which bind their second carrier by a few meV; that is the dominant athermal route once flashing has populated them. The phonon-ionization cross-sections follow from the golden rule for one-phonon absorption with the deformation-potential coupling of \cref{sec:6.2}, between effective-mass envelopes that include the valley--orbit splitting (six valleys in silicon, four in germanium), since the splitting sets the threshold. For $\mathrm{D}^-$ and $\mathrm{A}^+$ the bound state is a correlated two-electron envelope of Chandrasekhar type. Impact ionization by drifting carriers uses Born cross-sections on the same envelopes away from threshold and a distorted-wave calculation near it, where the hot tail of a Luke-limited distribution sits. The envelope basis is converged to \SI{1}{\percent} in binding energies and cross-sections, on a threshold grid no coarser than a fiftieth of the binding energy. A 216-atom supercell (\cref{app:E.10}), \SI{1.6}{nm} across, cannot hold a donor envelope of radius 2 to \SI{3}{nm} and supplies only the central-cell correction, at about $10^{3}$ core-hours with the shared \SI{60}{Ry} setting; the envelopes and the event-sequence driver take minutes.

Ramo's theorem stays exact with space charge, but the heat dissipated by drifting carriers differs from the induced-charge work by the change of electrostatic energy, which the energy balance of \cref{sec:19.5} carries; and flashing creates the $\mathrm{D}^-$ and $\mathrm{A}^+$ populations whose emission it was meant to remove.

\subsection{The mark-sufficiency test}

The two-stage factorization of \cref{sec:18.3,sec:19.6} lets the detector stage see the cascade only through a few marks (the pair number, the centroid and spread of the charge cloud, the phonon spectrum), and it is exact if the marks are a sufficient statistic for the seed class. The test compares the detector stage run on full cascade microstates with the same stage run on microstates resampled to match only the marks, for a dark-matter seed class and a photon seed class. The test statistic is the difference $\Delta P$ between the two arms in the distributions of the measured charge and of the phonon signal. A difference resolved at two standard deviations needs $2p(1 - p)(z/\Delta P)^2$ trajectories per arm, with $z = 2$, which for $\Delta P = 0.01$ at $p = 0.5$ is \num{2e4} per seed class. Common random numbers between the arms, as in \cref{app:E.12}, lower this, and the test costs $10^{2}$ to $10^{3}$ core-hours. The test can fail but cannot prove: a null difference can hold while the omitted features differ, if the detector stage happens to ignore them.

\subsection{Low-energy-excess inputs}

The collective-release model of \cref{sec:19.6} lacks two inputs: the relaxable elastic energy $U_{\mathrm{rel}}$ stored at cooldown, with its spread of release rates $[\varGamma_{\min}, \varGamma_{\max}]$ for a given holder and glue, and the branching ratio $b_{\mathrm{br}}$, the mean number of further releases triggered by the phonons of one release, which fixes the size law.

$U_{\mathrm{rel}}$ comes from a thermo-elastic finite-element calculation of the cooldown from \SI{300}{K} to \SI{10}{mK} for the crystal, the glue and the holder, with measured elastic tensors $C_{ij}(T)$ and expansion coefficients, or with the integrated thermal contraction of \cref{app:E.17} where no measurement exists. The elastic energy of each component is converged to \SI{5}{\percent} under a doubling of the mesh. The relaxable fraction and its spread in $\ln\varGamma$ come from low-temperature dynamic-mechanical data of the glue, not from the simulation. The energy law of single release events comes from molecular dynamics of a crystal--epoxy interface with $10^6$ atoms at 1 to \SI{10}{K}, which needs a machine-learned potential for cured epoxy (\cref{app:E.11}). None exists; one has to be trained on dispersion-corrected density-functional energies of the cured network and the interface, validated on elastic constants and thermal expansion, and extended to bond breaking by active learning. Event energies are read from energy drops, and the law is stable when the box grows from $10^6$ to \num{4e6} atoms. Molecular dynamics reaches nanoseconds, so it gives the energy law and never the rates, which come from the finite-element step.

The branching ratio needs the release-path catalogue of \cref{app:E.22}. The phonon pulse of one release is propagated by the transport Monte Carlo of \cref{app:E.12}, and a catalogued site with barrier $E_b$ triggers with probability $1 - \exp[-\int\dd t\,\varGamma(t)]$, where $\varGamma(t)$ is its phonon-assisted rate in the absorbed flux above $E_b$. The triggering volume is the integral of this probability over space, and $b_{\mathrm{br}}$ is the density of sites with barriers below the phonon energy times that volume, averaged by kinetic Monte Carlo over $10^6$ seeds. It is first-generation and linear in the flux, and it has to change by less than \SI{10}{\percent} when the catalogue grows from 50 to 100 paths, the phonon packets from $10^5$ to \num{4e5}, and the box from $L$ to $2L$. An incomplete catalogue, missing the low barriers, biases $b_{\mathrm{br}}$ low; a box smaller than the phonon mean free path truncates the triggering volume; and the thermal bath and the pulse must not both drive the same rate. The finite-element part costs $10^{2}$ to $10^{3}$ core-hours, the molecular dynamics $10^{3}$ to $10^{4}$ GPU-hours. A fit of the model to the binned rates of a glued and an unglued crystal~\cite{anthonypetersen2024stress} then tests the predicted scaling with $U_{\mathrm{rel}}$ without free parameters.

\subsection{Collection kernels of a stated detector}

The collection contrasts $\varpi_O$ and $\varpi_{\mathrm{gap}}$, the residual position spread and the collection ratio $\mathfrak{r}_{\mathrm{coll}}$, defined in \cref{sec:19.5,sec:20.5}, are linear functionals of the collection kernels $\mathfrak{C}_s(E_{\mathrm{ath}};\vect{r}_{\mathrm{ev}})$ of \cref{eq:20.5-kernel}. Electron and nuclear recoils enter only through their birth spectra, never through the kernels.

No kernel has been computed for a real detector. Two geometries are specified: the $1 \times 1 \times \SI{0.4}{cm^3}$ silicon crystal of Ref.~\cite{ren2021design}, whose top face carries two sensor arrays, and a germanium cylinder \SI{1}{cm} thick of the CDMS type. \Cref{tab:E25-runs} gives the grids: 40 logarithmic points in $E_{\mathrm{ath}}$ between 0.41 and \SI{62}{meV} (0.1 to \SI{15}{THz}), $8^3$ event positions in the symmetry-reduced octant, and $10^4$ phonons launched per cell as monochromatic sources, which is what makes the tally a kernel. Luke phonons are generated by the carrier transport at biases of 2, 4 and \SI{10}{V}. The approximation level, including the lumped films, is that of \cref{app:E.12}. Each geometry uses its actual sensor films; the mean-field estimate of \cref{sec:20.5}, by contrast, takes one film with a niobium-like gap ($2\Delta_{\mathrm{BCS}} = \SI{3.0}{meV}$) and one of aluminum ($2\Delta_{\mathrm{BCS}} = \SI{0.36}{meV}$).

The kernels are accepted on numerical and physical grounds. $\varpi_O$ must change by less than 0.01 when the energy grid is refined, and the statistical error per cell must fall below \SI{1}{\percent} when the launch count grows from $10^4$ to \num{4e4}. In the limit of diffuse, grey surfaces and diffusion faster than absorption, the position-averaged kernels must reproduce the mean-field solution of \cref{sec:20.5}, with a bare-surface loss of about $10^{-3}$ per bounce~\cite[Sect.~III\,G]{leman2012invited}; there $\varpi_O$ lies between 0.02 and 0.39 for mean Luke energies from 4.7 down to \SI{1.5}{meV} and the position spread is zero by construction, so the computed spread is the first new number the kernels give. Three systematic scans make up the nuisance of the contrasts: the surface coefficients varied between the acoustic- and diffuse-mismatch values, the probability of specular reflection, which no code default fixes, and, for germanium, the anharmonic constant over its factor-of-ten range.

Luke-phonon birth spectra depend on the field, so they must come from the carrier transport at each bias; a sampler that assigns one phonon energy per emission fixes $\mathfrak{r}_{\mathrm{coll}}$ by construction. Phonons below $2\Delta_{\mathrm{BCS}}$ of aluminum are lost energy, not absorbed energy, unless a tungsten sensor or another film takes them, and that assignment sets the sign of $\mathfrak{r}_{\mathrm{coll}} - 1$ (\cref{sec:20.5}); it must be stated with every result. Accelerated propagation modes that skip scattering events are checked against full tracking before use; and the transport code has to pass the film validation of \cref{app:E.25-film} first. The cost is about $2\times10^{12}$ propagation steps, some 140 core-hours per geometry; a first pass at $10^3$ phonons per cell takes under an hour at about \SI{3}{\percent} per cell.

\subsection{Gap-energy release and created-phonon spectra}

Each pair created at zero field carries the gap energy $\Eg$, which comes back as phonons when the pair recombines, and where and as what spectrum it comes back decides $\varpi_{\mathrm{gap}}$. Its phonon-borne part has a birth density $N_{\mathrm{eh}}\Eg\sum_\ell\pi_\ell(\vect{r},t)\,s_\ell(E_{\mathrm{ath}})$ summed over release channels $\ell$: bulk recombination of excitons by Auger or multiphonon processes, recombination at a surface, collection at an electrode, and quasiparticle injection into a film. Here $\pi_\ell$ is the channel flux in space and time and $s_\ell$ the spectrum emitted per channel. Then $\varpi_{\mathrm{gap}}$ is the sum over channels of the time-integrated channel fraction times the contrast of $s_\ell$ born at the release site, read from the kernels above. Bulk release of a hard spectrum, above the memory-loss scale of \cref{sec:20.2}, contributes nothing.

Besides the kinetic rates of \cref{app:E.24}, the release needs the exciton diffusion constant $D_X$, the surface recombination velocity $S_{\mathrm{surf}}$ and the spectra $s_\ell$. The channel fluxes come from a one-dimensional implicit finite-difference solution of exciton diffusion with Robin faces, with $\dd z \le L_X/20$ for the diffusion length $L_X = (D_X/\varGamma_X)^{1/2}$ and exciton decay rate $\varGamma_X$. An exciton born at depth $z_0$ recombines at a face with probability $[S_{\mathrm{surf}}/(S_{\mathrm{surf}} + D_X/L_X)]\,\ee^{-z_0/L_X}$, which reduces to $\ee^{-z_0/L_X}$ only for $S_{\mathrm{surf}} \to \infty$. The fluxes enter the transport as time-tagged sources; the released fraction is converged to $10^{-3}$ under grid and time step, with $S_{\mathrm{surf}}$ scanned from 10 to \SI{e4}{cm/s}. The surface and electrode spectra cannot be assumed: taking them as a hard line at the optical-phonon energy decides $\varpi_{\mathrm{gap}} = 0$ in advance. Until they are computed, they are scanned as mixtures $(1 - \kappa_s)s_{\mathrm{hard}} + \kappa_s s_{\mathrm{soft}}$, with $\kappa_s \in \{0, 0.25, 0.5, 0.75, 1\}$ and $s_{\mathrm{soft}}$ a Debye spectrum cut at 2, 4 or \SI{8}{meV}. Soft surface spectra are measured: in germanium at $\pm$\SI{3}{V}, recombination phonons from the surfaces carry \SI{13}{\percent} of the phonon signal in a low-frequency, ballistic regime~\cite[Sect.~V\,A]{leman2012invited}. The scan is replaced by the multiphonon line shape of the recombination center (a dangling bond at the Si/SiO$_2$ interface, or a state of a metal contact), with the Huang--Rhys factor and local modes from defect supercells (\cref{app:E.10}) converged to \SI{10}{\percent}, at about $10^3$ core-hours per center. At millikelvin $D_X$ and the binding rates depend on the impurity census of the actual crystal, droplets lie outside linear kinetics, and luminescence belongs to $E_{\mathrm{esc}}$, not to $\varpi_{\mathrm{gap}}$.

The primary phonons are treated the same way. The created-phonon spectra of electron and nuclear recoils at $E_R$ = \SIlist{50;100;300;1000}{eV}, branch-resolved on a \SI{0.5}{meV} grid with the zero-point energy of the final state subtracted (\cref{app:E.23}), are convolved with the kernels to give the arrival spectra and the observable difference between the two recoil types. \Cref{sec:20.2} bounds that difference by the $L^1$ distance of the normalized birth spectra below the memory-loss scale plus a term $\Delta_{\mathrm{h}}$ from the hard part; the bound reduces to the soft-fraction difference only when the spectral difference has one sign below that scale. The test is that refining the grid to \SI{0.25}{meV}, or moving the memory-loss scale by a factor of two, changes the bound by no more than $\Delta_{\mathrm{h}}$. Thermal phonons absorbed during the recoil enter as negative sources, never clipped, and the energy the molecular-dynamics thermostat removed enters with its own spectrum. The convolution costs nothing; a direct transport check samples births from the tables at about $10^8$ steps per seed set, minutes.

\subsection{Collection efficiencies of the calibration detectors}

The silicon phonon-yield ratio $Y_{\mathrm{ph}} = 0.952$ is degenerate between stored energy and collection (\cref{sec:19.5}). Breaking the degeneracy needs the collection efficiencies $f^{(p)}_{\mathrm{coll}}$ of each population $p$, primary phonons of nuclear and of electron recoils, Luke phonons and gap-energy phonons, for the detectors that measured the data. These are a CDMS~II silicon detector, \SI{76}{mm} across and \SI{10}{mm} thick at \SI{4}{V/cm}~\cite[Sect.~2]{agnese2018nuclear}, and a SuperCDMS germanium detector at $\pm$\SI{2}{V} with a $^{206}$Pb surface source~\cite{agnese2018energy}. The gap part of the collection offset $\delta_{\mathrm{coll}}$ is fixed by $\rho_g$, the efficiency for gap-energy phonons over that for hard primaries; $\rho_g - 1 = 0.21$ alone would account for the silicon deficit $1 - Y_{\mathrm{ph}} = 0.048$ beyond the 0.006 that the pair energy of the reconstruction produces, so $\rho_g$ must be known to $\pm 0.02$. The two-bias inversion of \cref{sec:19.5}, at 4 and \SI{40}{V}, assumes that $\mathfrak{r}_{\mathrm{coll}}$, the ionized fraction $f_{\mathrm{ion}}$ and the collection offset do not change with bias and that the reconstruction keeps its pair energy; the run computes the residual from the bias dependence of $\mathfrak{r}_{\mathrm{coll}}$, and a third bias tests all three.

The calculation adds charge transport to the kernels: Herring--Vogt valleys, Luke emission, intervalley scattering inside the field range of its fitted parameters, and trapping at constant cross-section (\cref{app:E.12}), with event positions from source maps of a Geant4 simulation~\cite{agostinelli2003geant4} of the neutron, gamma and lead sources. A carrier collected at an electrode releases its kinetic energy plus its energy relative to the Fermi level of that contact, supplied per contact; booking $\Eg/2$ per carrier with no phonons would set $\rho_g$ by construction. Each population needs $10^5$ events, and the inversion must be consistent between the two biases. The surface source of the germanium data makes the escape term of $\delta_{\mathrm{coll}}$ the leading one there. The cost is about $10^{11}$ steps, some 8 core-hours, per geometry.

\subsection{Anisotropic film interfaces}

The escape time $t_l = 4d_{\mathrm{f}}/(v_s\bar t_{\mathrm{esc}})$ of a phonon from a sensor film (\cref{sec:20.3}) needs the escape probability $\bar t_{\mathrm{esc}}$ of the actual interface. The isotropic acoustic- and diffuse-mismatch values differ by a factor 1.31 in $t_l$ for aluminum on silicon. The welded interface between two anisotropic elastic media, with Christoffel-matched plane waves and the six-by-six boundary problem for displacement and traction, gives the transmissions $t_\nu(\theta, \phi)$ per mode and incidence direction for Al(111) on Si(001) and for tungsten on silicon, in both directions. Thin-film interference enters through a transfer matrix when the wavelength is comparable to the thickness, and $t_l(E_{\mathrm{ath}})$ is tabulated for $d_{\mathrm{f}}$ from 40 to \SI{600}{nm}. The aluminum constants are the measured ones extrapolated to zero temperature, the set of \cref{sec:20.3}; finite-strain density-functional constants (\cref{app:E.4}), half an hour for fcc aluminum at a \SI{60}{Ry} cutoff, a $16^3$ grid and \SI{0.02}{Ry} smearing, serve only where no measurement exists, as for tungsten films.

The transmitted and reflected fluxes sum to one to $10^{-14}$; in the isotropic limit the code returns the escape probabilities of \cref{sec:20.3} for aluminum on silicon, $\bar t_{\mathrm{esc}} = 0.368$ (acoustic mismatch) and 0.282 (diffuse mismatch); the angle averages satisfy the reciprocity $\mathfrak{n}_1\bar t_{12} = \mathfrak{n}_2\bar t_{21}$ of \cref{sec:20.3}; and $\bar t$ is stable to $10^{-3}$ under refinement of the angular grid. The anisotropic result for Al(111) on Si(001) is compared with the sensitivity of \cref{sec:20.3}, where giving isotropic silicon its $[100]$ or $[111]$ speeds moves $\bar t_{\mathrm{esc}}$ to 0.327 or 0.405; a value outside that range points to the treatment of the aluminum side. The film formula carries its own hypotheses. Its flux average of the per-angle absorption is $1 - 2E_3(2d_{\mathrm{f}}/\mathfrak{l}_{\mathrm{f}})$, with $E_3$ the exponential integral, which the lumped form $1 - \ee^{-4d_{\mathrm{f}}/\mathfrak{l}_{\mathrm{f}}}$ of \cref{app:E.12} matches only to first order: with the escape probability set to one in both, the lumped form lies \SI{3}{\percent} to \SI{13}{\percent} above it for $d_{\mathrm{f}}/\mathfrak{l}_{\mathrm{f}}$ from 0.01 to 0.2. And $t_l$ holds for a film population that is isotropized faster than it escapes: without mixing in the film, transverse modes beyond \SI{37.7}{\degree} are totally reflected at Al/Si, some \SI{56}{\percent} of the film's phonon flux, and never escape. Real interfaces carry oxide or amorphous interlayers, and the fitted silicon-to-aluminum value of \cref{app:E.12} is an absorption probability $\bar t\,P_{\mathrm{abs}}$, not an escape probability. The whole calculation costs 2 to 10 core-hours.

\subsection{Validation of the film cascade}\label{app:E.25-film}

Every number that \cref{chap:20} takes from a lumped film rests on the quasiparticle yield $Y_{\mathrm{qp}}(E_{\mathrm{ath}})$ of \cref{sec:20.4}, and the code that produces it is tested before any kernel is trusted. Without film loss, $Y_{\mathrm{qp}}$ must equal $2\Delta_{\mathrm{BCS}}/E_{\mathrm{ath}}$ to round-off below $4\Delta_{\mathrm{BCS}}$, reach the maximum of 0.600 at $5.46\,\Delta_{\mathrm{BCS}}$ and settle at the 0.590 that the Kaplan cascade of \cref{sec:20.4} gives; the test points are $4.5\Delta_{\mathrm{BCS}}$, $6\Delta_{\mathrm{BCS}}$ and 10 to $100\Delta_{\mathrm{BCS}}$. With the escape law of a \SI{600}{nm} aluminum film, $Y_{\mathrm{qp}}$ falls with energy at 3, 5, 10 and $40\Delta_{\mathrm{BCS}}$, because the fastest phonons escape before they break pairs (\cref{sec:20.4}). Energy closes per event to $10^{-14}$ relative for incident phonons from 50 to $350\Delta_{\mathrm{BCS}}$. The last range matters because germanium prompt phonons near \SI{36}{meV} arrive at an aluminum film with some $200\Delta_{\mathrm{BCS}}$. The closure counts escaped energy as phonons actually returned to the substrate, so that an escape bin cannot absorb lost energy. The mean Luke-phonon energy against field is compared with the carrier model of \cref{app:E.12}. Each test runs in seconds and a benchmark of a few thousand events in minutes. Below $100\Delta_{\mathrm{BCS}}$ the yields agree with the loss-free forms within statistical error; above it the closure test decides.

\begin{table}[t]
  \caption{Event-level transport: grids, acceptance tests and cost of each calculation. Energies of phonons are $E_{\mathrm{ath}}$; core-hours are CPU core-hours unless marked GPU. Costs are per geometry or per parameter set as stated.}
  \label{tab:E25-runs}
  \centering
  \small
  \begin{tabular}{p{0.20\linewidth}p{0.29\linewidth}p{0.28\linewidth}p{0.11\linewidth}}
    \toprule
    quantity & grid & acceptance & cost \\
    \midrule
    Charge and exciton kinetics & $t_{\mathrm{win}}$: \SI{1}{\micro s} to \SI{1}{s}; thresholds 0.3--\SI{16}{eV} & $Y_{\mathrm{ph}} = 0.952$ (Si), \SI{6.08}{\percent} (Ge), Hecht limit & $\le 10^3$ core-h per set \\
    Athermal emission, space charge & envelope basis to \SI{1}{\percent}; $\dd z \le d/200$ & saturation of the field distortion & $\sim 10^3$ core-h \\
    Mark sufficiency & \num{2e4} trajectories per arm & $\Delta P$ resolved to 0.01 & $10^2$--$10^3$ core-h \\
    Low-energy-excess inputs & mesh $\times 2$; $10^6 \to \num{4e6}$ atoms; catalogue $50 \to 100$ & elastic energy per component to \SI{5}{\percent}; $b_{\mathrm{br}}$ to \SI{10}{\percent} & $10^2$--$10^3$ core-h; $10^3$--$10^4$ GPU-h \\
    Collection kernels & 40 energies, 0.41--\SI{62}{meV}; $8^3$ positions; $10^4$ phonons per cell & $\varpi_O$ to 0.01; mean-field limit & $\sim 140$ core-h \\
    Gap-energy release & $\dd z \le L_X/20$; $S_{\mathrm{surf}}$ 10--\SI{e4}{cm/s} & released fraction to $10^{-3}$ & seconds; $10^3$ per center \\
    Created-phonon spectra & \SI{0.5}{meV} $\to$ \SI{0.25}{meV} & bound stable to $\Delta_{\mathrm{h}}$ & minutes \\
    Collection efficiencies & $10^5$ events per population & $\rho_g$ to $\pm 0.02$; two biases & $\sim 8$ core-h \\
    Film interfaces & angular grid; $d_{\mathrm{f}}$ 40--\SI{600}{nm} & flux sum to $10^{-14}$; isotropic 0.368, 0.282; reciprocity & 2--10 core-h \\
    Film cascade validation & 3--$350\Delta_{\mathrm{BCS}}$ & $2\Delta_{\mathrm{BCS}}/E_{\mathrm{ath}}$; 0.590; closure $10^{-14}$ & $< 1$ core-h \\
    \bottomrule
  \end{tabular}
\end{table}
